\documentclass[instructions.tex]{subfiles}
\begin{document}
 
\chapter{De la communion.}
 
Communiez; communiez dignement; communiez souvent.
 
\primo Communiez. Jésus-Christ nous y invite tous\footnote{Venite ad me omnes. Matth.11.28}. Il le commande sous peine de mort : Si vous ne mangez la chair du Fils de l’homme, vous n’aurez point la vie en vous. Le corps ne peut vivre long-temps sans nourriture; comment votre âme pourrait-elle vivre long-temps de la vie de la grâce, et persévérer dans la sainteté, sans cette nourriture céleste? Combien de fidèles languiraient dans le service de Dieu, combien tomberaient dans de grands désordres, s’ils n’étaient pas soutenus par ce pain des forts! Combien seraient morts dans le trouble et le désespoir, s’ils n’avaient pas reçu le Dieu de toute consolation, et ce pain de vie dans le saint viatique!
 
\secundo Communiez dignement. S’il fallait tant de préparation pour participer aux mystères de l’ancienne loi, qui n’étaient que des figures et des ombres, quelle préparation ne faut-il pas pour recevoir un Dieu \footnote{Si in figurà tante praeparatiom quanta in veritate. S.Ambr.}! Saint Jean-Baptiste, le plus grand des prophètes, se reconnaissait indigne de se prosterner devant le Sauveur. Oh! combien plus devons-nous nous avouer indignes d’approcher de ce Dieu de sainteté!
 
Il faut, comme l’ordonne saint Paul, s’éprouver soi-même; s’exciter à des sentimens de foi, à une humble et sincère détestation de ses fautes, à un grand amour, à un grand désir de s’unir à son Dieu; sonder sa conscience, la purifier au moins de tous péchés mortels, avant de recevoir ce pain sacré. La nourriture n’est que pour les vivans; elle ne profite point aux morts : il faut donc être vivant, être en état de grâce, pour recevoir la divine Eucharistie. Communier en péché mortel, ou avec une affection volontaire au péché mortel, est un sacrilége énorme.
 
Oh! que cet attentat est injurieux à Jésus-Christ! Les Juifs le crucifièrent sans le connaître; mais l’indigne communiant le connaît, le reçoit et l’outrage! Ceux-là le crucifièrent sur le Calvaire, et celui-ci le crucifie dans son cœur, en unissant le Saint des saints avec le péché! union monstrueuse, dont le Sauveur a infiniment plus d’aversion, que du supplice de la croix.
 
Un tyran fit autrefois attacher des hommes vivans à des cadavres infects, pour les suffoquer. On fait quelque chose de semblable envers le Sauveur, lorsque, par une sacrilége communion, on l’unit à une âme souillée de crimes, dont il a plus d’horreur que la personne la plus délicate n’en a de l’infection d’un cadavre.
 
S’en prendre à la personne de Jésus-Christ par l’indigne communion, oh! que ce crime est funeste! Loin d’y recevoir la vie, l’âme y reçoit la mort\footnote{Mors est malis. Prose.}. En recevant son juge, on reçoit son jugement, dit saint Paul \footnote{Judicium sibi manducat. Cor.11} ; c’est-à-dire que Jésus-Christ prononce la condamnation, et scelle de son sang l’arrêt de mort du téméraire qui le reçoit indignement.
 
Ecoutons saint Chrysostome sur le sort infortuné de Judas. Judas est murmurateur, Jésus-Christ le souffre; il est avare et larron, Jésus-Christ le souffre; il fait un complot pour trahir son maître, Jésus-Christ le souffre : mais communie-t-il indignement, il est livré sur-le-champ à la puissance de Satan \footnote{Post buccellam intriovit in eum Satanas. S. Chrys.}. L’indigne communion attire de grands châtimens. Saint Cyprien en rapporte plusieurs exemples tragiques; et saint Chrysostome reprochait à son peuple, que la profanation des saints mystères était la cause des calamités et des fléaux qui désolaient l’empire.
 
Le roi Sédécias donne un soufflet à un prophète; cet attentat méritait une punition, cependant Dieu le souffre. Mais le lévite Oza porte-t-il sa main sur l’arche sainte, cinquante-deux mille Bethsamites osent-ils la dévoiler, ils sont tous sur-le-champ punis de mort; le roi Ozias, tout saint qu’il était, ne fut-il pas puni de Dieu, pour avoir mis la main à l’encensoir? et le roi Balthazar massacré pour avoir profané les vases du temple? O mon Dieu! que doivent donc craindre ceux qui profanent aujourd’hui ce qu’il y a de plus saint dans la religion, et qui s’en prennent par le sacrilége à la personne même de Jésus-Christ! Quel compte n’en rendront pas les ministres de Dieu, si, par leur lâcheté, ils coopèrent à ces profanations!
 
Si vous êtes dans le sacrilége ou dans l’habitude du crime, c’est pour vous un grand malheur; mais ce malheur n’est pas sans remède. Allez vous prosterner aux pieds de Jésus-Christ; pleurez l’affront que vous lui avez fait; et, quoi qu’il vous coûte, faites une confession sincère qui vous convertisse; quelque grand pécheur que vous soyez, il aura pitié de vous. De même qu’autrefois il mangeait plus volontiers à la table des humbles pécheurs qu’à celle des superbes pharisiens; aussi n’est-il personne qui soit plus favorablement reçu de Jésus-Christ qu’un pécheur touché de repentir, qui se convertit pour s’unir à son Dieu.
 
\tertio Communiez souvent. On ne doit pas s’éloigner de la communion, lorsque, par une humble confession, on est rentré dans la grâce de Dieu. Jésus-Christ a donné ce sacrement sous le symbole du pain, pour nous apprendre qu’il devait être la nourriture ordinaire de notre âme. L’Eglise nous invite, nous exhorte et nous conjure d’y participer souvent, à l’exemple des premiers fidèles.
 
Notre plus grand regret devrait être de nous voir privés de cette nourriture céleste. Nos faiblesses et nos dangers doivent nous faire sentir le besoin que nous avons de ce pain quotidien. Plus on s’en éloigne, plus on est faible dans le service de Dieu. Peut-on vivre saintement, en s’éloignant du sacrement institué pour nous soutenir dans la sainteté?
 
C’est donc une vaine excuse inspirée par l’illusion ou par la lâcheté, de dire qu’on ne communie pas souvent, pour mieux s’y disposer, ou pour ne pas abuser du sacrement. Plus vous communierez souvent avec foi, plus dignement vous le ferez. Ce n’est pas en faisant rarement une chose, qu’on apprend à la bien faire.
 
Au reste, la communion fréquente n’est pas la sainteté ni la perfection, mais elle est le moyen qui nous y conduit. La santé ne consiste pas à manger souvent, mais à profiter de la nourriture; de même la vie sainte ne consiste pas précisément à communier souvent, mais à profiter de la communion, et à se corriger de plus en plus. C’est pourquoi saint Augustin ayant dit : Prenez chaque jour ce pain qui peut vous profiter chaque jour , ajoute : Vivez de telle sorte que vous méritiez de le recevoir tous les jours.
 
\section{Résolutions}
 
\primo Il faut que je m’efforce de si bien vivre désormais, que mon confesseur puisse me permettre de communier souvent.
 
\secundo Mais je ne m’approcherai jamais de la table sainte, le péché mortel dans le cœur.
 
\tertio Et je me préparerai toujours à cette grande action, en excitant en moi de vifs sentimens de foi, d’anéantissement, de confiance, de désir, de contrition et d’amour.
 
\prayer{O Jésus! que votre tendresse pour nous est grande! et qui pourra jamais assez l’admirer, assez y correspondre?}
 
\end{document}
